\section{Application Scenario \& Protocol Interaction}
\subsection{Application scenario}
The application is a command-line Hybrid FTP system that separates the
control and data planes. FTP commands, standard three-digit replies,
authentication, and session state use a persistent TCP connection. File
contents and directory listings use UDP with a custom application-layer
Go-Back-N reliability protocol.

The server listens at \code{127.0.0.1:2121}. Each accepted TCP connection
creates one independent \code{Session} in one daemon thread. The session owns
the authenticated user, working directory, transfer type, transfer mode,
active/passive endpoint state, rename state, and cancellation event. Server
files are confined to \code{server_files/}; client downloads are written to
\code{client_downloads/}.

\subsection{Architecture}

\begin{figure}[H]
\centering
\begin{tikzpicture}[
  node distance=8mm and 14mm,
  box/.style={draw,rounded corners,align=center,minimum width=45mm,minimum height=9mm},
  arrow/.style={-{Latex[length=2.5mm]},thick}
]
\node[box] (user) {User commands\\login, put, get, ls, verify};
\node[box,below=of user] (cli) {\code{client/client.py}\\Interactive CLI};
\node[box,below=of cli] (fc) {\code{client/ftp_client.py}\\Client protocol API};
\node[box,right=45mm of fc] (server) {\code{server/server.py}\\TCP accept + thread/session};
\node[box,below=of server] (session) {\code{server/session.py}\\Authentication and dispatch};
\node[box,below=18mm of fc,xshift=30mm] (rdt) {\code{common/rdt.py}\\Go-Back-N reliable UDP};
\draw[arrow] (user)--(cli);
\draw[arrow] (cli)--(fc);
\draw[arrow,<->] (fc)--node[above]{TCP commands/replies}(server);
\draw[arrow] (server)--(session);
\draw[arrow,<->] (fc)--node[left]{UDP data/ACK}(rdt);
\draw[arrow,<->] (rdt)--(session);
\end{tikzpicture}
\caption{Project-wide control and data paths.}
\end{figure}

\subsection{Full TCP + UDP lifecycle}

\begin{longtable}{p{0.07\textwidth}p{0.18\textwidth}p{0.18\textwidth}p{0.47\textwidth}}
\toprule
\# & Sender & Receiver & Interaction\\
\midrule
1 & Client & Server & TCP \code{connect()} to port 2121; server replies \code{220}.\\
2 & Client & Server & \code{USER admin}; server replies \code{331}.\\
3 & Client & Server & \code{PASS ********}; server verifies the selected user and replies \code{230}.\\
4 & Client & Server & \code{PASV}; server binds UDP port 50000--51000 and replies \code{227}.\\
5 & Client & Server & \code{STOR remote.bin} over TCP and an UDP \code{RDY} datagram to the PASV endpoint.\\
6 & Server & Client & \code{150 Opening data connection}.\\
7 & Client & Server & Window of UDP DATA packets containing sequence, length, flags, CRC-32 and payload.\\
8 & Server & Client & Cumulative ACK for the highest contiguous sequence. Loss causes timeout/NAK and GBN retransmission.\\
9 & Server & File system & Reassemble, validate representation, and write under \code{server_files/}.\\
10 & Server & Client & \code{226 Transfer complete}.\\
11 & Client & Server & \code{HASH SHA256 remote.bin}; server replies with \code{213} and digest.\\
12 & Client & Client & Calculate local digest and compare end to end.\\
13 & Client & Server & \code{QUIT}; server replies \code{221}, closes sockets, and deregisters the session.\\
\bottomrule
\end{longtable}

\subsection{Passive and active modes}

\textbf{PASV.} The server binds a random UDP socket, sends its endpoint in
\code{227}, and learns the client's UDP endpoint from the client's RDY
datagram. PASV is the default.

\textbf{PORT.} The client binds a UDP socket and advertises it with
\code{PORT h1,h2,h3,h4,p1,p2}. The server stores that endpoint and sends RDY
from its per-transfer UDP socket. The client learns the server endpoint from
that datagram.

\subsection{Approved command coverage}

\begin{tabularx}{\textwidth}{>{\bfseries}p{0.28\textwidth}X}
\toprule
Category & Commands\\
\midrule
Authentication/session & USER, PASS, QUIT, NOOP, HELP\\
Navigation/directories & PWD, CWD, CDUP, MKD, RMD\\
Listing/metadata & LIST, NLST, STAT, SIZE, MDTM\\
Representation/data setup & TYPE, MODE, PORT, PASV\\
File transfer & RETR, STOR, STOU, APPE, ABOR\\
Management/integrity & DELE, RNFR, RNTO, HASH\\
\bottomrule
\end{tabularx}

All 28 commands required by the specification have server dispatch handlers.

\subsection{Evaluation-level mapping}

\begin{longtable}{p{0.15\textwidth}p{0.32\textwidth}p{0.43\textwidth}}
\toprule
Level & Requirement & Implementation\\
\midrule
Basic & Authentication & Stateful USER/PASS on the TCP session\\
Basic & ASCII and upload/download & TYPE A, STOR and RETR\\
Basic & Fixed data mechanism & PASV default\\
Advanced & Binary files & TYPE I, byte-preserving payload\\
Advanced & Directory tree & Navigation, listing, create/remove, metadata\\
Advanced & Active/passive & Runtime PORT/PASV switching\\
Advanced & Concurrency & One thread and Session per TCP client\\
Excellent & Reliable UDP & Sequence, ACK/NAK, CRC, timeout, retransmission\\
Excellent & Flow/congestion control & GBN window, Slow Start, AIMD\\
Excellent & End-to-end integrity & MD5/SHA-256 local/remote comparison\\
\bottomrule
\end{longtable}
